\documentclass[11pt]{article}
\usepackage{campaign}

\title{Draft Legislation: An Act to Amend the Criminal Code (Medical Assistance in Dying)}
\author{Justin Bogner}
\date{2026}

\begin{document}

\maketitle

\begin{abstract}
\noindent \textit{Note to Parliamentarians and Legislative Counsel: This draft legislation is designed to amend Section 241.2 of the Criminal Code. It resolves the ongoing deadlock regarding Track 2 MAiD and MD-SUMC (Mental Disorder as Sole Underlying Medical Condition) by replacing the subjective ``grievous and irremediable medical condition'' standard with an objective ``decisional capacity'' standard.}
\end{abstract}

\section*{1. SHORT TITLE}
This Act may be cited as the \textit{Capacity-Based Medical Assistance in Dying Act}.

\section*{2. AMENDMENTS TO THE CRIMINAL CODE}

\textbf{(1) Section 241.2 (1) of the Criminal Code is replaced by the following:}

\vspace{0.5em}
\noindent \textbf{Eligibility for medical assistance in dying} \\
\textbf{241.2 (1)} A person may receive medical assistance in dying only if they meet all of the following criteria:
\begin{enumerate}
    \item[(a)] they are eligible --- or, but for any applicable minimum period of residence or waiting period, would be eligible --- for health services funded by a government in Canada;
    \item[(b)] they are at least 18 years of age;
    \item[(c)] they possess durable decisional capacity, meaning the ability to understand and appreciate the nature and consequences of the request for medical assistance in dying;
    \item[(d)] they have made a voluntary request for medical assistance in dying that, in particular, was not made as a result of external pressure; and
    \item[(e)] they give informed consent to receive medical assistance in dying.
\end{enumerate}

\textbf{(2) Section 241.2 (2) of the Criminal Code (Grievous and irremediable medical condition) is repealed.}

\vspace{1em}
\noindent \textbf{(3) Section 241.2 (3.1) of the Criminal Code (Safeguards — natural death not reasonably foreseeable) is replaced by the following:}

\vspace{0.5em}
\noindent \textbf{Safeguards — two-stage capacity assessment} \\
\textbf{(3.1)} Before a medical practitioner or nurse practitioner provides medical assistance in dying to a person, the practitioner must:
\begin{enumerate}
    \item[(a)] be of the opinion that the person meets all of the criteria set out in subsection (1);
    \item[(b)] ensure that the person’s request for medical assistance in dying was made in writing and signed and dated by the person;
    \item[(c)] be satisfied that the request was signed and dated after the person received the initial capacity assessment (Stage One);
    \item[(d)] ensure that there are at least 90 clear days between the day on which the request was signed and the day on which medical assistance in dying is provided where the person's natural death is not reasonably foreseeable; where natural death is reasonably foreseeable, no waiting period is required and provision may occur as soon as is feasibly possible;
    \item[(e)] ensure that a second, independent assessment (Stage Two) is conducted by a medical practitioner with specialized training in psychiatry;
    \item[(f)] be satisfied that the independent psychiatric evaluator has confirmed that the person possesses durable decisional capacity; and
    \item[(g)] if the person has a mental illness or psychiatric disorder, be satisfied by the independent psychiatric evaluator that the illness or disorder does not actively impair the person's ability to make a rational, informed, and durable decision.
\end{enumerate}

\section*{3. AMENDMENTS REGARDING DATA COLLECTION AND SOCIOECONOMIC REPORTING}

\textbf{(1) Section 241.31 (1) of the Criminal Code is amended by adding the following:}

\vspace{0.5em}
\noindent \textbf{Information to be provided — Socioeconomic Drivers (The structural accountability mechanism)} \\
\textbf{(1.1)} Any medical practitioner or nurse practitioner who receives a written request for medical assistance in dying must conduct a standardized assessment of the socioeconomic factors driving the request, and must provide the designated recipient with the following information:
\begin{enumerate}
    \item[(a)] whether the person reported an inability to access or afford necessary medical, palliative, or psychiatric care;
    \item[(b)] whether the person reported inadequate housing or impending homelessness; and
    \item[(c)] whether the person reported insufficient disability support services.
\end{enumerate}

\noindent \textbf{(1.2) Non-disqualification} \\
Under no circumstances shall the presence of the socioeconomic drivers listed in subsection (1.1) invalidate the eligibility of a person who has been determined to possess durable decisional capacity under Section 241.2 (1). The information is collected solely for the purpose of health system accountability and parliamentary oversight.

\end{document}
